\section{Flujo m\'aximo (Ford--Fulkerson)}
\whenuse{Red peque\~na con capacidades enteras y flujo total bajo, o como
implementaci\'on conceptual r\'apida. Para casos grandes preferir Dinic.}

Busca caminos aumentantes con DFS y actualiza el grafo residual. Esta versi\'on
usa capacidades enteras no negativas.

Cada arista representa cu\'anto flujo puede enviarse desde \texttt{source} hasta
\texttt{sink}. Al enviar $x$, se resta $x$ de la capacidad directa y se suma a
la inversa; esta arista residual permite deshacer o redirigir flujo anterior.
Ford--Fulkerson repite: buscar un camino con capacidad, enviar el m\'inimo de sus
capacidades y terminar cuando ya no existe otro camino.

\begin{lstlisting}[style=cpp]
struct FordFulkerson {
    struct Edge { int to, reverse; ll capacity; };
    int n;
    vector<vector<Edge>> graph;
    vector<char> seen;

    explicit FordFulkerson(int size)
        : n(size), graph(size), seen(size) {}

    void add_edge(int from, int to, ll capacity) {
        assert(capacity >= 0);
        Edge forward{to, (int)graph[to].size(), capacity};
        Edge backward{from, (int)graph[from].size(), 0};
        graph[from].push_back(forward);
        graph[to].push_back(backward);
    }

    ll dfs(int vertex, int sink, ll pushed) {
        if (vertex == sink) return pushed;
        seen[vertex] = true;
        for (Edge &edge : graph[vertex]) {
            if (seen[edge.to] || edge.capacity == 0) continue;
            ll flow = dfs(edge.to, sink, min(pushed, edge.capacity));
            if (flow == 0) continue;
            edge.capacity -= flow;
            graph[edge.to][edge.reverse].capacity += flow;
            return flow;
        }
        return 0;
    }

    ll max_flow(int source, int sink) {
        assert(source != sink);
        ll answer = 0;
        while (true) {
            fill(seen.begin(), seen.end(), false);
            ll pushed = dfs(source, sink, LINF);
            if (pushed == 0) break;
            answer += pushed;
        }
        return answer;
    }
};
\end{lstlisting}

\complexity{$O(EF)$ para capacidades enteras, donde $F$ es el flujo m\'aximo;
es pseudopolinomial; memoria $O(V+E)$. Con capacidades grandes suele convenir Dinic.}

\section{Flujo m\'aximo (Dinic)}
\whenuse{Capacidades, cortes m\'inimos, asignaciones con cupos y selecci\'on con
restricciones que se modelan como una red fuente--sumidero.}

Calcula el mismo flujo m\'aximo que Ford--Fulkerson, pero agrupa los aumentos.
Un BFS construye niveles usando aristas residuales y los DFS env\'ian todo el
flujo posible avanzando s\'olo al nivel siguiente. \texttt{next\_edge} evita
revisar aristas que ya se demostraron in\'utiles en la fase actual. Es la opci\'on
usual de flujo m\'aximo en competencia.

\begin{lstlisting}[style=cpp]
struct Dinic {
    struct Edge { int to, reverse; ll capacity; };
    int n;
    vector<vector<Edge>> graph;
    vector<int> level, next_edge;

    explicit Dinic(int size)
        : n(size), graph(size), level(size), next_edge(size) {}

    void add_edge(int from, int to, ll capacity) {
        Edge forward{to, (int)graph[to].size(), capacity};
        Edge backward{from, (int)graph[from].size(), 0};
        graph[from].push_back(forward);
        graph[to].push_back(backward);
    }

    bool bfs(int source, int sink) {
        fill(level.begin(), level.end(), -1);
        queue<int> q;
        level[source] = 0; q.push(source);
        while (!q.empty()) {
            int v = q.front(); q.pop();
            for (const Edge &edge : graph[v])
                if (edge.capacity > 0 && level[edge.to] == -1)
                    level[edge.to] = level[v] + 1, q.push(edge.to);
        }
        return level[sink] != -1;
    }

    ll dfs(int v, int sink, ll pushed) {
        if (v == sink || pushed == 0) return pushed;
        for (int &i = next_edge[v]; i < (int)graph[v].size(); ++i) {
            Edge &edge = graph[v][i];
            if (edge.capacity == 0 || level[edge.to] != level[v] + 1) continue;
            ll flow = dfs(edge.to, sink, min(pushed, edge.capacity));
            if (flow == 0) continue;
            edge.capacity -= flow;
            graph[edge.to][edge.reverse].capacity += flow;
            return flow;
        }
        return 0;
    }

    ll max_flow(int source, int sink) {
        assert(source != sink);
        ll answer = 0;
        while (bfs(source, sink)) {
            fill(next_edge.begin(), next_edge.end(), 0);
            while (ll pushed = dfs(source, sink, LINF)) answer += pushed;
        }
        return answer;
    }
};
\end{lstlisting}

\complexity{general $O(V^2E)$ tiempo; memoria $O(V+E)$; suele ser mucho m\'as
r\'apido en pr\'actica.}

\section{Flujo de costo m\'inimo}
\whenuse{Enviar cierta cantidad de flujo optimizando un costo: transporte,
matching ponderado con capacidades o asignaci\'on de recursos.}

SPFA por aumento; admite costos negativos si no hay ciclos negativos de costo
alcanzables con capacidad residual positiva.

Cada arista tiene capacidad y costo por unidad. En cada iteraci\'on se busca el
camino residual de menor costo, se env\'ia todo lo permitido y se acumula
\texttt{flujo $\times$ costo}. \texttt{run(s,t,k)} intenta enviar hasta $k$
unidades y devuelve \texttt{\{flujo\_enviado,costo\_total\}}; si no existe
capacidad suficiente, el flujo devuelto es menor que $k$.

\begin{lstlisting}[style=cpp]
struct MinCostMaxFlow {
    struct Edge { int to, reverse; ll capacity, cost; };
    int n;
    vector<vector<Edge>> graph;

    explicit MinCostMaxFlow(int size) : n(size), graph(size) {}

    void add_edge(int from, int to, ll capacity, ll cost) {
        Edge forward{to, (int)graph[to].size(), capacity, cost};
        Edge backward{from, (int)graph[from].size(), 0, -cost};
        graph[from].push_back(forward);
        graph[to].push_back(backward);
    }

    pair<ll,ll> run(int source, int sink, ll flow_limit = LINF) {
        ll total_flow = 0, total_cost = 0;
        vector<ll> distance(n);
        vector<int> parent_vertex(n), parent_edge(n);
        while (total_flow < flow_limit) {
            fill(distance.begin(), distance.end(), LINF);
            vector<char> in_queue(n, false);
            queue<int> q;
            distance[source] = 0; q.push(source); in_queue[source] = true;
            while (!q.empty()) {
                int v = q.front(); q.pop(); in_queue[v] = false;
                for (int i = 0; i < (int)graph[v].size(); ++i) {
                    Edge const &edge = graph[v][i];
                    if (edge.capacity == 0 ||
                        distance[edge.to] <= distance[v] + edge.cost) continue;
                    distance[edge.to] = distance[v] + edge.cost;
                    parent_vertex[edge.to] = v; parent_edge[edge.to] = i;
                    if (!in_queue[edge.to]) q.push(edge.to), in_queue[edge.to] = true;
                }
            }
            if (distance[sink] == LINF) break;
            ll pushed = flow_limit - total_flow;
            for (int v = sink; v != source; v = parent_vertex[v])
                pushed = min(pushed,
                    graph[parent_vertex[v]][parent_edge[v]].capacity);
            for (int v = sink; v != source; v = parent_vertex[v]) {
                Edge &edge = graph[parent_vertex[v]][parent_edge[v]];
                edge.capacity -= pushed;
                graph[v][edge.reverse].capacity += pushed;
            }
            total_flow += pushed;
            total_cost += pushed * distance[sink];
        }
        return {total_flow, total_cost};
    }
};
\end{lstlisting}

\complexity{$O(FVE)$ en el peor caso para flujo entero $F$; memoria $O(V+E)$.}

\subsection{Elecci\'on r\'apida}

\begin{itemize}
    \item S\'olo maximizar cantidad: preferir \textbf{Dinic}.
    \item Maximizar cantidad y minimizar costo: \textbf{Min-Cost Max-Flow}.
    \item Ford--Fulkerson: implementaci\'on conceptual o flujo total peque\~no.
\end{itemize}

En matching bipartito se modela
\texttt{source $\to$ izquierda $\to$ derecha $\to$ sink}, normalmente con
capacidades $1$. En la versi\'on ponderada, el costo de la arista central es el
costo de realizar esa asignaci\'on.
